\chapter{Sore Throat and Headache}

\section*{Definition}
Sore throat with headache commonly accompanies viral respiratory infection but may occur with streptococcal infection, influenza, COVID-19, sinus disease, or meningitis.

\section*{When to See a Doctor}
Seek urgent evaluation for Stiff neck, confusion, breathing or swallowing difficulty, drooling, severe unilateral swelling, dehydration, rash, or rapidly worsening illness. Arrange routine or prompt assessment for persistent, recurrent, unexplained, or function-limiting symptoms.

\section*{How It Develops}
Sore throat with headache commonly accompanies a viral upper-respiratory infection, but streptococcal pharyngitis, influenza, COVID-19, sinus disease, and less common deep-neck or meningeal infection must be considered from the examination and associated symptoms.

\section*{Causes}
\begin{itemize}
  \item Viral infection, streptococcal pharyngitis, influenza, COVID-19, sinusitis, dehydration, and less commonly meningitis or deep-neck infection.
  \item Medication effects, environmental exposures, and chronic disease may contribute.
  \item Less common but important causes include malignancy, vascular disease, or organ failure when symptoms persist or progress.
\end{itemize}

\section*{Related Symptoms}
\begin{itemize}
  \item Pain, fever, fatigue, nausea, or changes in normal organ function
  \item Bleeding, weight change, shortness of breath, weakness, or neurologic symptoms when a serious cause is present
\end{itemize}

\section*{Medical Evaluation}
\begin{itemize}
  \item History addresses onset, duration, severity, triggers, exposures, medications, medical conditions, pregnancy status when relevant, and warning symptoms.
  \item Examination focuses on vital signs and the organ systems suggested by the symptom.
  \item Testing may include blood or urine studies, cultures, imaging, physiologic testing, or specialist examination according to the clinical pattern.
\end{itemize}

\section*{Treatment Options}
Supportive care and age-appropriate analgesia are used for most viral illness. Antibiotics are reserved for confirmed or strongly suspected bacterial disease, while severe swallowing difficulty, drooling, neck stiffness, or breathing problems require urgent care.

\section*{Self-Care and Home Management}
Avoid known triggers, maintain reasonable hydration, record timing and associated symptoms, and use only treatments appropriate for the suspected cause. Do not use leftover antibiotics or delay urgent care because symptoms temporarily improve.

\section*{References}
\begin{thebibliography}{99}
\bibitem{sore_throat_and_headache:idsa} Shulman ST, Bisno AL, Clegg HW, et al. Clinical practice guideline for the diagnosis and management of group A streptococcal pharyngitis: 2012 update by the Infectious Diseases Society of America. \textit{Clin Infect Dis}. 2012;55:e86--e102.
\bibitem{sore_throat_and_headache:escmid} Pelucchi C, Grigoryan L, Galeone C, et al. Guideline for the management of acute sore throat. \textit{Clin Microbiol Infect}. 2012;18(suppl 1):1--28.
\bibitem{sore_throat_and_headache:harrison} Jameson JL, et al. \textit{Harrison's Principles of Internal Medicine}. 21st ed. McGraw-Hill; 2022.
\end{thebibliography}
